\section{The Claim}

\textbf{[STIPULATION]} The gravitational constant $G$ is in this theory not an
independent constant of nature but a unit bridge -- like $c$, $\hbar$, and $k_B$
(A080).

From this follows the stronger and better-known claim of the corpus: the problem of
quantum gravity does not arise in this form. This document examines what supports
this claim and what does not.

\section{Why $G$ Is a Bridge}

\textbf{[B]} $G$ connects a mass with a length. In the natural system of the
theory, in which $\tilde T\cdot m = 1$ holds and time and mass are the same
quantity in two readings, this connection does not need to be separately
established -- it is already in the duality.

This is the same statement as for $\hbar$: there the bridge between mass and time
is meant, here the one between mass and length. Both are, in a system that does not
carry time, mass, and length as independent units, conversions and not natural laws.

\textbf{[STIPULATION]} Operationally $G$ remains indispensable. One who computes
in SI needs it. Ontologically it is secondary. This is exactly the dual role that
A080 describes for all bridge constants, and it is not a demotion.

\section{Why the QG Problem Does Not Arise in This Form}

\textbf{[B]} The usual problem runs: quantise gravitation as a field and thereby
obtain a renormalisable theory. It arises because the geometry is treated as a
dynamic field that must in turn be quantised.

In this theory the geometry is \emph{pre-ordered} and not dynamic in the same
sense: the $T^4/Z_3$ is stipulated (A020), the quantum structure arises as its
spectrum (A030, A070). There is therefore nothing left that would need to be
subsequently quantised.

\textbf{Honest assessment of this statement.} It is not a solved problem but a
\emph{dissolved} one -- and dissolved by a stipulation. One who does not stipulate
the pre-ordered geometry still has the problem. The theory does not answer the
question; it does not pose it.

This is the \emph{negative} argument. A positive argument stands alongside it,
which derives the numerical value of $G$ from $\xi$: $G=\xi^2/(4\,m_\text{char})$
(A145). Both belong together -- the negative explains why quantisation is
unnecessary, the positive explains where the value comes from.

This is a legitimate move, but it must not be presented as a solution. The
difference is the same as between a derivation and a calibration (A130): both are
permissible, but they are not the same.

\section{What Follows from This and What Does Not}

\textbf{[B]} It follows: no gravitons as necessary constituents, no
renormalisation problems of gravitation, no singularities from a field quantisation.

\textbf{[STIPULATION]} It does \emph{not} follow: that the classical gravitational
dynamics is reproduced. The A-series shows nowhere that the observed gravitational
interaction with its strength and inverse-square law follows from the compact
geometry. This is a major open point and is not minimised here.

\section{Verification Script}

\begin{itemize}
  \item The bridge role of $G$: that $G$ can be absorbed in natural units just
    like $c$, $\hbar$, and $k_B$, shown by the consistency of the unit equations:
    \texttt{python/A\_Serie\_Skripte/a140\_g\_bruecke.py}
\end{itemize}

\section{Open Questions}

\textbf{First, gravitational dynamics has not yet been treated as an A-series
document -- it is taken up in A142.} The complete field-equation framework
(time-field Lagrange, scalar/fermion coupling, modified Dirac equation, Higgs
coupling, $\xi$ from Higgs matching) is available in the legacy corpus and is
incorporated in A142. What remains after that: (a) the complete quantisation of
these field equations, (b) the transition to measurable predictions beyond the
Schwarzschild case, (c) the connection to the projection chain (A090).

\textbf{Second, the strength of the interaction is not derived.} Treating $G$ as a
bridge does not explain why gravitation is forty orders of magnitude weaker than
electromagnetism.

\textbf{Third, the dissolution of the QG problem is a consequence of the
stipulation.} See above; booked as such, not as a result.

\section{Supersedes}

This document subsumes: Dok.~012 (gravitational constant), Dok.~127 (gravitation),
Dok.~163 (quantum gravity). Incorporated: P40.

\section{Use of AI Tools}

This document was created with the assistance of AI tools. The questions,
stipulations, and all substantive decisions come from the author; the tools
formulate, compute, and create verification scripts.
Responsibility for content and errors rests entirely with the author.
Full wording of the rules in A010.
